\documentclass{report}

\usepackage[T2A]{fontenc}
\usepackage[utf8]{inputenc}
\usepackage{pgfplots}


\usepackage[sort,compress]{cite}
\usepackage{amsmath}
\usepackage{amssymb}
\usepackage{amsthm}
\usepackage{fancyvrb}
\usepackage{longtable}
\usepackage{array}
\usepackage[english,russian]{babel}


%добавленные мной
\usepackage{geometry}
\geometry{papersize={27.4 cm,25.4 cm}}
\geometry{left=1cm}
\geometry{right=1cm}
\geometry{top=2cm}
\geometry{bottom=2cm}

\usepackage{listings}
\usepackage{color}

\definecolor{mygray}{rgb}{0.4,0.4,0.4}
\definecolor{mygreen}{rgb}{0,0,15.6}
\definecolor{myorange}{rgb}{1.0,0.4,0}

\lstset{
basicstyle=\footnotesize\sffamily\color{black},
commentstyle=\color{mygray},
frame=single,
numbers=left,
numbersep=5pt,
numberstyle=\tiny\color{mygray},
keywordstyle=\color{mygreen},
showspaces=false,
showstringspaces=false,
stringstyle=\color{myorange},
tabsize=2
}
%конец добавления


\begin{document}

\title{Задание 1, Ассемблер}
\author{Якунина Владимира Александровича}
\maketitle

\section{Текст задания.}

\begin{enumerate}
\item Измените программы из примеров 1, 2 и 3 так, чтобы они выводили на экран ваши фамилию, имя и номер группы. Используя командные файлы, подготовьте к выполнению и запустите программы из примеров 1, 2 и 3. Убедитесь, что они выводят на экран нужный текст и успешно завершаются.
\item Заполните таблицы трассировки для программ вывода строки на экран из примеров 1, 2 и 3.
\end{enumerate}

\section{Тексты трех программ на языке ассемблера с комментариями.}

\lstinputlisting[language=C++]{ex1.m}

\lstinputlisting[language=C++]{ex2.m}

\lstinputlisting[language=C++]{ex3.m}

\section{Тексты командных файлов.}

\lstinputlisting[language=C++]{m1.b}

\lstinputlisting[language=C++]{m2.b}

\lstinputlisting[language=C++]{m3.b}

\section{Таблицы трассировки программ.}

 \begin{table}[h!]
	\begin{tabular}{|c|c|c|c|c|c|c|c|c|c|c|c|c|}
	\hline
	Шаг & Машинный код & Команда & AX & BX & CX & DX & SP & DS & SS & CS & IP & CZSOPAID \\
	\hline
1 & B8BD48 & MOV AX,48BD & 0000 & 0000 & 0000 & 0000 & 0100 & 489D & 48AD & 48BF & 0000 & 00000010 \\
  	\hline
2 & 8ED8 & MOV DS,AX & 48BD & 0000 & 0000 & 0000 & 0100 & 489D & 48AD & 48BF & 0003 & 00000010 \\
  	\hline
3 & B409 & MOV AH,09 & 48BD & 0000 & 0000 & 0000 & 0100 & 48BD & 48AD & 48BF & 0005 & 00000010 \\
  	\hline
4 & BA0000 & MOV DX,0000 & 09BD & 0000 & 0000 & 0000 & 0100 & 48BD & 48AD & 48BF & 0007 & 00000010 \\
  	\hline
5 & CD21 & INT 21 & 0000 & 09BD & 0000 & 0000 & 0100 & 48BD & 48AD & 48BF & 000A & 00000010 \\
  	\hline
 6 & B8004C & MOV AX,4C00 & 09BD & 0000 & 0000 & 0000 & 0100 & 48BD & 48AD & 48BF & 000C & 00000010 \\
  	\hline
7 & CD21 & INT 21 & 0000 & 4C00 & 0000 & 0000 & 0100 & 48BD & 48AD & 48BF & 000F & 00000010 \\
  	\hline
  	\end{tabular}
	\caption{ТРАССИРОВКА `EX1'}\label{tab1}
\end{table}

 \begin{table}[h!]
	\begin{tabular}{|c|c|c|c|c|c|c|c|c|c|c|c|c|}
	\hline
	Шаг & Машинный код & Команда & AX & BX & CX & DX & SP & DS & SS & CS & IP & CZSOPAID \\
	\hline
1 & B8AF48 & MOV 48,AF & 0000 & 0000 & 0000 & 0000 & 0100 & 489D & 48B1 & 48AD & 0000 & 00000010 \\
  	\hline
2 & 8ED8 & MOV DS,AX & 48AF & 0000 & 0000 & 0000 & 0100 & 489D & 48B1 & 48AD & 0003 & 00000010 \\
  	\hline
3 & B409 & MOV AH,09 & 48AF & 0000 & 0000 & 0000 & 0100 & 48AF & 48B1 & 48AD & 0005 & 00000010 \\
  	\hline
4 & BA0000 & MOV DX,0000 & 09AF & 0000 & 0000 & 0000 & 0100 & 48AF & 48B1 & 48AD & 0007 & 00000010 \\
  	\hline
5 & CD21 & INT 21 & 0000 & 09AF & 0000 & 0000 & 0100 & 48AF & 48B1 & 48AD & 000A & 00000010 \\
  	\hline
 6 & B8004C & MOV AX,4C00 & 09AF & 0000 & 0000 & 0000 & 0100 & 48AF & 48B1 & 48AD & 000C & 00000010 \\
  	\hline
7 & CD21 & INT 21 & 0000 & 4C00 & 0000 & 0000 & 0100 & 48AF & 48B1 & 48AD & 000F & 00000010 \\
  	\hline
  	\end{tabular}
	\caption{ТРАССИРОВКА `EX2'}\label{tab2}
\end{table}


 \begin{table}[h!]
	\begin{tabular}{|c|c|c|c|c|c|c|c|c|c|c|c|c|c|c|c|}
	\hline
Шаг & Машинный код & Команда & AX & BX & CX & DX & SP & DS & SS & CS & IP & CZSOPAID & SI & DI & BP \\
	\hline
1 & 59 & POP CX & 0000 & 0000 & 0000 & 0000 & 0000 & 489D & 48AC & 48AD & 0100 & 00000010 & & & \\
  	\hline
2 & 61 & POPA &  &  &  &  & 0002 &  &  &  & 0101 & & & & \\
  	\hline
3 & 6B756E69 & IMUL SI,[DI+6E],00 &  & 6864 & 0048 & 4400 & 0012 &  &  &  & 0102 & & 6568 & 6474 & 706C \\
  	\hline
4 & 6E & OUTSB &  &  &  &  &  &  &  &  & 0106 & & 0000 &  &  \\
  	\hline
5 & 20566C & AND [BP+6C],DC &  &  &  &  &  &  &  &  & 0107 & & 0001 &  &  \\
  	\hline
6 & 61 & POPA &  &  &  &  &  &  &  &  & 010A & 01001010 &  &  &  \\
  	\hline
7 & 64696D697220 & IMUL BP,FS:[DI+69] &  & 0000 & 0000 & 0000 & 0022 &  &  &  & 010B &  & 0000 & 0000 & 0000 \\
  	\hline  	
8 & 3335 & XOR SI,[DI] &  &  &  &  &  &  &  &  & 0111 & 11011010 &  &  & 4804 \\
  	\hline  	
9 & 3120 & XOR [BX+SI],SP &  &  &  &  &  &  &  &  & 0113 & 00000010 & 20CD &  &  \\
  	\hline  
10 & 43 & INC BX &  &  &  &  &  &  &  &  & 0115 &  &  &  &  \\
  	\hline   	
 11 & 53 & PUSH BX &  & 0001 &  &  &  &  &  &  & 0116 &  &  &  &  \\
  	\hline   	
 12 & 49 & DEC CX &  &  &  &  & 0020 &  &  &  & 0117 &  &  &  &  \\
  	\hline   	
  13 & 54 & PUSH SP &  &  & FFFF &  &  &  &  &  & 0118 & 00101110 &  &  &  \\
  	\hline   	
14 & 2124 & AND [SI],SP &  &  &  &  & 001E &  &  &  & 0119 & 00101110 &  &  &  \\
  	\hline  
15 & B409 & MOV AH,09 &  &  &  &  &  &  &  &  & 011B & 01001010 &  &  &  \\
  	\hline  
 16 & BA0001 & MOV DX,0100 & 0900 &  &  &  &  &  &  &  & 011D &  &  &  &  \\
  	\hline  
17 & CD21 & INT 21 &  &  &  & 0100 &  &  &  &  & 0120 &  &  &  &  \\
  	\hline  
18 & B8004C & MOV AX,4C00 &  &  &  &  &  &  &  &  & 0122 &  &  &  &  \\
  	\hline  
 19 & CD21 & INT 21 & 4C00 &  &  &  &  &  &  &  & 0125 &  &  &  &  \\
  	\hline   	
  	\end{tabular}
	\caption{ТРАССИРОВКА `EX3'}\label{tab3}
\end{table}

\section{Ответы на контрольные вопросы.}

1.  Запустите под управлением отладчика программу из примера 1. Как располагаются сегменты программы в памяти? Сопоставьте полученные данные с соответствующими размерами сегментов из листинга трансляции.

Ответ: CS – сегмент команд, DS – сегмент данных, SS – сегмент стека. CS = 48BEh, DS = 489Dh, SS = 48ADh

2.  Что такое сегментный (базовый) адрес?

Ответ: Адрес ячейки памяти. Находится по формуле: база сегмента + смещение.

3.  Почему перед началом выполнения программы из примера 1 содержимое регистра DS в точности на 10h меньше содержимого регистра SS?

Ответ: Адрес начала сегмента всегда кратен 16-ти (каждый сегмент должен начинаться на границе параграфа).

4.  Почему в программе из примера 1 первая команда mov AX, data, начинающаяся с байта 0000 сегмента команд, занимает 3 байта? Какая еще команда занимает 3 байта?

Ответ: Так как здесь используется прямая адресация. Ещё 3 байта занимает команда mov dx, 0000.

5.  Чему равен сегментный адрес сегмента команд программы из примера 2? Почему?

Ответ: 48ADh, так как используется модель памяти small.

6.  Сравните содержимое регистра SP в программах из примера 2 и примера 3. Объясните, как получены эти значения.

Ответ: Во 2-м примере SP = 0100h, так как используется модель памяти small и стек располагается в первых 256 байтах сегмента памяти, а в 3-м – FFFEh, так как используется модель памяти tiny и сегменты кода, данных и стека находятся в одном сегменте размером 64 Кб.

7.  Какие операторы называют директивами ассемблера? Приведите примеры директив.

Ответ: Директивы – это параметры, указывающие программе ассемблеру, каким образом следует объединять инструкции для создания модуля, который станет работающей программой. Примеры: ASSUME, SEGMENT.

8.  Зачем в последнем предложении end указывают метку, помечающую первую команду программы?

Ответ: Для того, чтобы программа ассемблер поняла какая из подпрограмм закончена.

9.  Как числа размером в слово хранятся в памяти и как они заносятся в 2-ух байтовые регистры?

Ответ: Они хранятся в памяти как шестнадцатеричное число и записываются справа налево. При занесении их в 2-х байтовые регистры байты меняются местами.

10.Как инициализируются в программе выводимые на экран текстовые строки?

Ответ: Сначала строка заносится в память также как и обычное число, а затем строка выводится на экран.

11.Что нужно сделать, чтобы обратиться к DOS для вывода строки на экран? Как DOS определит, где строка закончилась?

Ответ: Использовать оператор int 21h. В конец строки всегда добавляется символ \$.

12.Программы, которые должны исполняться как. EXE и. COM, имеют существенные различия по: размеру, сегментной структуре, механизму инициализации. Охарактеризуйте каждый из перечисленных пунктов.

Ответ: COM программы занимают меньше места в памяти и быстрее загружаются. В COM программах префикс программного сегмента, код программы, инициализированные данные и стек располагаются в одном сегменте. При запуске такой программы в начале строится PSP, занимающий 256 байт, далее располагается содержимое программы. Указатель стека устанавливается на конец сегмента.
EXE программы содержат сегменты кода, данных и стека. EXE файл загружается, начиная с адреса PSP:0100h. В процессе загрузки считывается информация EXE-заголовка в начале файла, при помощи которого загрузчик настраивает ссылки на сегменты в загруженном модуле, так как программа загружается в произвольный сегмент. После настройки ссылок управление передается загрузочному модулю к адресу CS:IP, извлеченному из заголовка EXE.



\end{document}